\subsection{Requerimientos Funcionales:Módulo 2}
 \begin{enumerate}
        \begin{enumerate}
            \item \textbf{Verificaciones de validez de las entradas}
                \begin{enumerate}
                    \item El sistema deberá verificar que los datos recibidos de la base de datos estén completos y no contengan valores nulos.
                    \item El sistema deberá verificar que los datos de pH estén dentro del rango esperado (4.0 a 7.0).
                    \item El sistema deberá verificar que los datos de temperatura estén dentro del rango esperado (0°C a 40°C).
                    \item El sistema deberá verificar que los datos de CO estén dentro del rango esperado (0 a 1000 ppm).
                    El sistema deberá verificar que los datos incluyan un timestamp válido y estén ordenados cronológicamente para garantizar la correcta secuencia temporal del análisis.
                \end{enumerate}   
            
            \item \textbf{Secuencia exacta de operaciones}
                \begin{enumerate}
                    \item El sistema deberá iniciar el análisis de datos automáticamente cada hora.
                    \item El sistema deberá procesar los datos en el orden en que fueron recibidos.
                    \item El sistema deberá generar un informe de análisis después de procesar un lote completo de datos.
                \end{enumerate}   
        
            \item \textbf{Respuestas a situaciones anormales}
                \begin{enumerate}
                    \item El sistema deberá generar una alerta si se detectan valores fuera del rango esperado en cualquiera de los parámetros (pH, temperatura, CO).
                    \item El sistema deberá registrar en un log cualquier error encontrado durante el análisis de datos.
                    \item El sistema deberá notificar al administrador si se detectan patrones anormales en los datos, como una caída rápida del pH o un aumento inusual en los niveles de CO.
                \end{enumerate}   
       
            \item \textbf{Efecto de los parámetros}
                \begin{enumerate}
                    \item El sistema deberá ajustar la frecuencia del análisis de datos si se detectan cambios significativos en los parámetros monitoreados.
                    \item El sistema deberá priorizar el análisis de datos si se detectan valores críticos en cualquiera de los parámetros.
                \end{enumerate}   
        
            \item \textbf{Relación entre las salidas y las entradas}
                \begin{enumerate}
                    \item El sistema deberá generar informes detallados basados en los datos analizados, incluyendo gráficos y tendencias.
                    \item El sistema deberá actualizar la base de datos con los resultados del análisis para su posterior consulta.
               \end{enumerate}
        \end{enumerate}
    \item \textbf{Módulo 3:Base de Datos}
        \begin{enumerate}
            \item \textbf{Verificaciones de validez de las entradas}
                \begin{enumerate}
                    \item El sistema deberá verificar que los datos recibidos del sistema de sensores estén completos y no contengan valores nulos.
                    \item El sistema deberá verificar que los datos de pH estén dentro del rango esperado (4.0 a 7.0).
                    \item El sistema deberá verificar que los datos de temperatura estén dentro del rango esperado (0°C a 40°C).
                    \item El sistema deberá verificar que los datos de CO estén dentro del rango esperado (0 a 1000 ppm).
                \end{enumerate}   
            
            \item \textbf{Secuencia exacta de operaciones}
                \begin{enumerate}
                    \item El sistema deberá almacenar los datos recibidos en la base de datos en el orden en que fueron recibidos.
                    \item El sistema deberá realizar copias de seguridad automáticas de la base de datos cada 24 horas.
                    \item El sistema deberá permitir consultas y recuperaciones de datos por parte del módulo de análisis de datos.
                \end{enumerate}   
        
            \item \textbf{Respuestas a situaciones anormales}
                \begin{enumerate}
                    \item El sistema deberá generar una alerta si se detectan intentos no autorizados de acceso a la base de datos.
                    \item El sistema deberá registrar en un log cualquier error encontrado durante las operaciones de almacenamiento o recuperación de datos.
                    \item El sistema deberá notificar al administrador si se detectan problemas recurrentes con la integridad de los datos.
                \end{enumerate}   
       
            \item \textbf{Efecto de los parámetros}
                \begin{enumerate}
                    \item El sistema deberá ajustar la frecuencia de las copias de seguridad si se detectan cambios significativos en el volumen de datos almacenados.
                    \item El sistema deberá optimizar las consultas si se detectan patrones comunes en las solicitudes de datos.
                \end{enumerate}   
        
            \item \textbf{Relación entre las salidas y las entradas}
                \begin{enumerate}
                    \item El sistema deberá proporcionar acceso rápido y eficiente a los datos almacenados para su análisis y visualización.
                    \item El sistema deberá mantener la integridad y seguridad de los datos almacenados en la base de datos.
                \end{enumerate}
        \end{enumerate}

    \item \textbf{Módulo 4:Página web} 
        \begin{enumerate}
            \item \textbf{Verificaciones de validez de las entradas}
                \begin{enumerate}
                    \item El sistema deberá verificar que los datos recibidos del módulo de análisis de datos estén completos y no contengan valores nulos.
                    \item El sistema deberá verificar que los datos de pH, temperatura y CO estén dentro del rango esperado antes de mostrarlos en la interfaz web.
                \end{enumerate}   
            
            \item \textbf{Secuencia exacta de operaciones}
                \begin{enumerate}
                    \item El sistema deberá actualizar la interfaz web con los datos más recientes cada 10 minutos.
                    \item El sistema deberá permitir a los usuarios filtrar y ordenar los datos mostrados en la interfaz web.
                    \item El sistema deberá proporcionar gráficos y visualizaciones basadas en los datos analizados.
                \end{enumerate}   
        
            \item \textbf{Respuestas a situaciones anormales}
                \begin{enumerate}
                    \item El sistema deberá mostrar un mensaje de error si no se pueden cargar los datos debido a problemas de conexión con la base de datos.
                    \item El sistema deberá registrar en un log cualquier error encontrado durante la actualización o visualización de datos en la interfaz web.
                    \item El sistema deberá notificar al administrador si se detectan problemas recurrentes con la carga de datos en la interfaz web.
                \end{enumerate}   
       
            \item \textbf{Efecto de los parámetros}
                \begin{enumerate}
                    \item El sistema deberá ajustar la frecuencia de actualización de la interfaz web si se detectan cambios significativos en los parámetros monitoreados.
                    \item El sistema deberá optimizar la visualización de datos si se detectan patrones comunes en las solicitudes de los usuarios.
                \end{enumerate}   
        
            \item \textbf{Relación entre las salidas y las entradas}
                \begin{enumerate}
                    \item El sistema deberá proporcionar una interfaz intuitiva y fácil de usar para que los usuarios puedan acceder y comprender los datos monitoreados.
                    \item El sistema deberá garantizar que los datos mostrados en la interfaz web sean precisos y reflejen el estado actual del monitoreo.
                \end{enumerate}
        \end{enumerate}

\end{enumerate}